\documentclass[11pt]{article}
\renewcommand{\familydefault}{\sfdefault}
\usepackage[usenames]{color}
\usepackage{graphicx}
\usepackage{psfig}
\usepackage{epsf}
\usepackage{html}
\usepackage{lscape}
\textheight 9.0in
\textwidth 6.0in
\hoffset -0.7in
\voffset -0.5in
\newcommand{\link}[1]{{\color{dblue}\htmladdnormallink{#1}\par}}
\newcommand{\m}[1]{\mbox{#1}}
\newcommand{\change}[1]{{\color{red}\it #1}}
\newsavebox{\fmbox}
\newenvironment{fmpage}
     {\begin{lrbox}{\fmbox}\begin{minipage}{6.5in}}
     {\end{minipage}\end{lrbox}\colorbox{iblue}{\fbox{\usebox{\fmbox}}}}

\newenvironment{fmppage}
     {\begin{lrbox}{\fmbox}\begin{minipage}{6.5in}}
     {\end{minipage}\end{lrbox}\colorbox{ipink}{\fbox{\usebox{\fmbox}}}}


\begin{document}

\vskip -1in
\psfig{file=ivoa.ps,width=5.0in}

\definecolor{ipink}{rgb}{1.0,0.937,0.957}
\definecolor{dblue}{rgb}{0.60,0.60,1.0}
\definecolor{ddblue}{rgb}{0.20,0.20,1.0}
\definecolor{iblue}{rgb}{0.9,0.9,1.0}
\huge
\color{dblue}
\vskip 0.2in
\par\noindent{\bf IVOA Spectral Data Model: SED Supplement}
\Large
\par\noindent{\bf Version 0.96}
\vskip 0.1in
\par\noindent{\bf IVOA Working Draft \\ 2005-09-18}
\color{Black}
\vskip 0.2in

\normalsize

\par\noindent{\bf This version:}

{\color{dblue}
\htmladdnormallink{http://www.ivoa.net/Documents/WD/SpectrumDM/SpectrumDM-2005XXXX.html}

}

\par\noindent{\bf Latest version:}

{\color{dblue}
\link{http://www.ivoa.net/Documents/latest/SpectrumDM.html}

}
\par\noindent{\bf Previous versions:}

{\color{dblue}
\link{http://hea-www.harvard.edu/$\sim$jcm/vo/docs/spec0.93.html}

}

\centerline{\it Note: ivoa.net links are not yet active; the above
are placeholders.}

\vskip 0.2in

\noindent{\bf Editors:}

Jonathan McDowell, Doug Tody

\noindent{\bf Authors:}

Jonathan McDowell, Doug Tody, Tamas Budavari,  Markus Dolensky, Frank Valdes,
Pavlos Protopapas, Arnold Rots,
and the IVOA Data Access Layer Working Group.



\begin{abstract}

We present a data model describing the aggregation of spectrophotometric
datasets with spectral and temporal coordinates and associated metadata
into a single SED (Spectral Energy Distribution).

\end{abstract}


\subsection*{ Status of this document }

This is a Working Draft, developed with the intention to support the
Simple Spectral Access Protocol. This is the first global release of the
document. The working group seeks confirmation that comments have been
addressed to the satisfaction of the community.

Comments on this document are due 30 Aug 2005 for consideration in the
next version of this document. They should be sent to the dal@ivoa.net
mailing list.

\par\noindent
This is an IVOA Working Draft for review by IVOA members and other
interested parties. It is a draft document and may be updated, replaced,
or obsoleted by other documents at any time. It is inappropriate to use
IVOA Working Drafts as reference materials or to cite them as other than
"work in progress." A list of current IVOA Recommendations and other
technical documents can be found at
\link{http://www.ivoa.net/Documents}.

This document has been developed with support from the National Science
Foundation's \link{http://www.nsf.gov}/ Information Technology Research
Program under Cooperative Agreement AST0122449 with The Johns Hopkins
University, from the UK Particle Physics and Astronomy Research Council
(PPARC) \link{http://www.pparc.ac.uk}, 
and from the Eurpean Commission's
Sixth Framework Program \link{http://fp6.cordis.lu/fp6/home.cfm}
via the
Optical Infrared Coordination Network (OPTICON),
\link{http://www.astro-opticon.org}.

The {\bf Virtual Observatory (VO)} is general term for a collection of
federated resources that can be used to conduct astronomical research,
education, and outreach. 

The {\bf International Virtual Observatory Alliance
(IVOA)} (\link{http://www.ivoa.net}) is a global collaboration of separately
funded projects to develop standards and infrastructure that enable VO
applications.

\tableofcontents

\newpage

\section{Introduction and Motivation}

One of the advantages of the Virtual Observatory is the ability
to intelligently combine data from different sources and then 
to further manipulate the combined data. In order to support
interoperable implementations of this capability, we need standard
ways of describing aggregated data that have the power to retain
the full metadata of their components. In this document we
describe the spectral version of such a dataset, the Spectral Energy
Distribution (SED). 
SED consist of multiple spectra, time series and photometry points, usually for a single
object. Our model is layered on the related IVOA Spectral data model.


\section{SED Data Model summary}

Our model for an SED is a set of spectra and/or time series, some of
which may have only one or few data points (photometry) and  each of
which may have different contextual metadata  (aperture, position,
etc.). An SED is made up by aggregating individual spectra/time-series
objects as 'segments'.

A general SED may be considered as consisting of {\bf segments}
for each of which the associated contextual metadata is constant. 

\begin{figure}[h]
\psfig{file=spec1.eps,width=5.0in}
\caption[a]{A moderately complicated SED with five segments a through e.
segments {\bf a} and {\bf b} are photometry points with associated
errors (vertical bars) and bandpass widths (horizontal bars).
segment {\bf c} is a time series with four measurements at
different times with the same instrumental configuration. segment {\bf d }
is an upper limit measurement. segment {\bf e } is a simple spectrum, with
point-to-point statistical errors indicated by vertical bars and
an overall (correlated) systematic error indicated by the dashed lines.}
\end{figure}

\begin{figure}[h]

\colorbox{iblue}{
\psfig{file=specuml6s.eps,width=6.0in}
}
{
\colorbox{iblue}{
\begin{minipage}{6.0in}
Figure 2: 
UML class diagram for the SED data model. The SED model
consists of a number of Segments each of which is consistent
with the Spectrum model, and which may represent a Spectrum,
Time Series, or photometry point. In addition, it may have
a Target object. A single SED object is assumed to correspond 
to data for a single
`target', i.e. the same target for each segment.  
\end{minipage}
}
}
\end{figure}


\section{Spectral data model Measurement objects}

\subsection{Spectrum and Time Series}



A simple Spectrum is an SED object for which the spectral
coordinate varies but the time coordinate
is fixed, and for which the associated metadata are constant 
(i.e. a single segment with a fixed time.)

A simple Time Series is an SED object for which the spectral
coordinate is constant but the time coordinate
varies, and for which the associated metadata are constant (i.e. one
set of metadata for the whole time series; a single segment
at a fixed spectral coordinate).

\subsection{SED attributes}

The overall SED object may contain values indicating the
number of SED segments and curation information about their
assembly into a single SED, as well as their overall
spectral range (the union of the segment bandpass coverages).
 
We introduce an SED.Type field with values {\bf observed }  or {\bf simulated }  or
{\bf composite } 
(default = observed).
The values indicate whether the data consists of
actual measurements versus theoretical or empirical calculations, or is a composite
of several observed datasets.
 


 {\bf OPTIONAL: The SED attributes are all optional except for DataModel, which
should have the value "SED-0.96". The NSegments
and SpectralMin/MaxWavelength can be derived from the metadata in the segments.
Note that the latter fields, if present, are required to be wavelength in meters
and are consistent with the Coverage.SpectralMinimumWavelength/SpectralMaximumWavelength
in the IVOA Resource Metadata document.
 }  

  

\begin{flushleft}
\colorbox{iblue}{\small
\begin{minipage}[l]{6.5in}
\begin{tabular}{lp{1.4in}p{1.6in}ll}
\multicolumn{3}{c}{SED Fields}\\

Field      & UCD1+  & Meaning & Unit  & Req \\
\hline 
 & & \\
 
SED.Creator         &      meta.id    & Person or organization   & & O\\
 
  & & creating the SED \\
 
SED.DocURL     &    meta.url           & URL for dataset documentation & & O\\
 
SED.Date           & time;meta.dataset      & Data processing/creation date& & O \\
 
SED.Type           &       & SED type, see above & & O\\
  
SED.NSegments       & meta.number          & Number of segments & & O \\
 
SED.SpectralMinimumWavelength  &em.wl;stat.min         & Total spectral coord range, wavelength& m & O\\
 
SED.SpectralMaximumWavelength   &em.wl;stat.max          & Total spectral coord range  & m & O\\

 
SED.DataModel    & meta.id    & IVOA Data Model version & &R\\
 
\end{tabular}
 \end{minipage}
}
\end{flushleft}

\begin{figure}[h]

\colorbox{iblue}{
\psfig{file=specu4b.eps,height=3.0in}
}
{
\colorbox{iblue}{
\begin{minipage}{2.0in}
Attributes for SED and Target objects
\end{minipage}
}
}

\end{figure}

\section{Serializations}

\subsection{FITS serialization}

The SED FITS serialization is a generalization of the single-segment
Spectrum serialization.

We define a reference serialization of this data model as a FITS binary table.
The table represents a single segement - spectrum or photometry point 
or time series - as a single row of a table.
Variable-length arrays may be used to contain the array quantities.
In each case below where a `variable length array' is specified,
fixed length arrays are
also acceptable if all rows of the table have the same length.
(The format is similar in spirit to the X-ray PHA type II spectral dataset).

We adopt the convention that columns which are constant (same value for all rows)
may if desired be omitted and the value given as a keyword instead. (e.g. the column
TTYPEn='INSTRUME'  is replaced by a keyword  INSTRUME = 'value'). This is a
small overhead in the FITS reading interface.

In particular, the SPCO\_UCD and FLUX\_UCD keywords give the spectral coordinate and flux
UCDs for the spectrum. 
In an SED if not all the segments use the same flux and spectral coordinate
UCDs these keywords become columns.

If each segment has the same spectral coordinate UCD, the spectral coordinate 
may also be identified by optional 1Sn\_1 and 1CTYPn keywords
as described in the Spectrum serialization.

Several extra keywords are defined for a SED dataset:

\vskip 0.2in
\colorbox{ipink}{\small
\begin{minipage}[l]{5.0in}
\begin{tabular}{lll} 
FITS keyword  & Data model field   & Value \\
\hline
 & & \\
 
VOCREATE      & SED.Creator \\
 
VOCRID        & SED.CreatorID \\
 
DATE          & SED.Date \\
 
BANDPASS      & SED.Bandpass \\

 
\end{tabular}
\end{minipage}
}
\vskip 0.1in

In the Spectrum model, the main data columns were declared
as fixed-length array columns.
Since different SED segments may have different lengths,
FITS' variable-length array mechanism may instead be used for these columns
to save space. However, fixed-length arrays may also be used; in this
case the SIZE column must be used to 
specify the number of points in each segment.


We summarize this with a full sample FITS extension header with
multiple SED segments, using variable-length arrays.

{ \color{ddblue} \footnotesize
\begin{verbatim}
XTENSION= 'BINTABLE'           / binary table extension
BITPIX  =                    8 / 8-bit bytes
NAXIS   =                    2 / 2-dimensional binary table
NAXIS1  =                   80 / width of table in bytes
NAXIS2  =                    6 / number of rows in table
PCOUNT  =                  208 / size of special data area
GCOUNT  =                    1 / one data group (required keyword)
TFIELDS =                   15 / number of fields in each row
EXTNAME = 'SPECTRUM '           / name of this binary table extension
VOCLASS = 'Spectrum V0.96'      / VO Data Model
VOCREATE= 'MMT Archive'       / VO Creator
VOCRID  = 'ivoa://cfa.harvard.edu' / VO Publisher ID URI
DATE    = '2004-08-30T14:18:17' / Date and time of file creation
DATE-OBS= '2004-06-03T21:18:17' / Date and time of observation
BANDPASS= ' '    / SED.Bandpass
RA_NOM       =     233.72789197     / [deg] Nominal RA    
DEC_NOM      =      23.49792615     / [deg] Nominal Dec   
OBJECT  = 'ARP 220 '           / Source name
SRCCLASS= 'Galaxy'             /
SPECTYPE= 'ULIRG'              /
RA_TARG      =     233.73791700     / Observer's specified target RA
DEC_TARG     =      23.50333300     / Observer's specified target Dec
REDSHIFT=              0.01812  / Emission redshift
DS_IDENT= 'cfa://whatever'   /
TITLE   = 'Observations of Merging Galaxies' /
VERSION =  2                     / Reprocessed 2004 Aug
VOPUB   = 'CfA Archive'        / VO Publisher authority
VOLOGO  = 'http://cfa.harvard.edu/vo/cfalogo.jpg' / VO Creator logo
SPCO_UCD= 'em.wl'                    /
FLUX_UCD= 'phot.fluDens;em.wl'     /
EQUINOX =  2.0000000000000E+03 / default
RADECSYS= 'ICRS    '           / default
TIMESYS = 'TT     '           / Time system
TIMEUNIT= 's'                / Time unit                
SPECSYS = 'TOPOCENT'            / Wavelengths are as observed
TIMESDIM= 'T'                  / Time SIDim 
SPECSDIM= '10-10 L'             / Spectral SIDim
FLUXSDIM= '10+7 ML-1T-3'       / Flux SDim
MJDREF  =  0.0                 / MJD zero point for times
SPEC_RES=              5.0  / [Angstrom] Spectral resolution
SKY_RES =              1.0  / [arcsec] Spatial.Resolution

TELESCOP= 'MMT '           / Telescope
INSTRUME= 'BCS    '           / Instrument
FILTER  = 'G220    '           / Grating
COMMENT  ---------------------------
COMMENT  WCS Paper 3 Keywords
1S11_1   = 'SPCO'            / Column name with spectral coord
1CTYP11  = 'WAVE-TAB'        / Spectral coord is WAVE 
1S12_1   = 'SPCO'            / Column name with spectral coord
1CTYP12  = 'WAVE-TAB'        / Spectral coord is WAVE 
1S13_1   = 'SPCO'            / Column name with spectral coord
1CTYP13  = 'WAVE-TAB'        / Spectral coord is WAVE 
1S14_1   = 'SPCO'            / Column name with spectral coord
1CTYP14  = 'WAVE-TAB'        / Spectral coord is WAVE 
COMMENT  ---------------------------
TTYPE1  = 'INSTRUME  '           / Instrument ID
TFORM1  = '8A      '           / format of field
TTYPE2  = 'FILTER  '           / Filter ID
TFORM2  = '8A      '           / Format
TTYPE3  = 'RA '                / Position RA of aperture center
TFORM3  = '1D      '           /
TUNIT3  = 'deg     '           /
TTYPE4  = 'DEC     '                / Position Dec of aperture center
TFORM4  = '1D      '           /
TUNIT4  = 'deg     '           /
TTYPE5  = 'APERTURE'           / Aperture diameter (physical or extraction)
TFORM5  = '1E       ' /
TUNIT5 =  'arcsec '   /
TTYPE6 = 'TIME'
TFORM6 = '1D '
TUNIT6 =  'd'  /  MJD days
TTYPE7 = 'EXPOSURE' /  Effective exposure time
TFORM7 = '1E'
TUNIT7 = 's' 
TTYPE8 = 'SYS_ERR' / Fractional systematic error
TFORM8 = '1E'
TUNIT8 = ' '
TTYPE9 = 'SPCO'   / Wavelength
TFORM9 = '1PE'
TUNIT9 = 'Angstrom'
TTYPE10= 'SPCO_LO' /
TFORM10= '1PE'
TUNIT10= 'Angstrom'
TTYPE11= 'SPCO_HI' /
TFORM11= '1PE'
TUNIT11= 'Angstrom'
TTYPE12= 'FLUX' /
TFORM12= '1PE'
TUNIT12= 'erg cm**(-2) s**(-1) Angstrom**(-1)'
TTYPE13= 'ERR_LO' /
TFORM13 = '1PE'
TUNIT13 ='erg cm**(-2) s**(-1) Angstrom**(-1)'
TTYPE14= 'ERR_HI' /
TFORM14 = '1PE'
TUNIT14 ='erg cm**(-2) s**(-1) Angstrom**(-1)'
TTYPE15= 'QUALITY' /
TFORM15 = '1PI'
\end{verbatim}
}

The data would look like
{\tiny
\begin{verbatim}
MMT/BCS  G300  233.73791  23.50333 2.0 52984.301203 1500.0 0.15  3200.0 3195.0 3205.0 1.48E-12 2.0E-14 2.0E-14   0
                                                                 3210.0 3205.0 3215.0 1.52E-12 3.0E-14 3.0E-14   0
                                                                 3220.0 3215.0 3225.0 0.38E-12 0.38E-12 0.0      0
                                                                 3230.0 3225.0 3235.0 1.62E-12 3.0E-14 3.0E-14   0
MMT/BCS  G300  233.73792  23.50334 2.0 52102.103211 1480.0 0.15  3200.0 3195.0 3205.0 3.48E-12 2.0E-14 2.0E-14   0
                                                                 3210.0 3205.0 3215.0 2.52E-12 3.0E-14 3.0E-14   0
                                                                 3220.0 3215.0 3225.0 1.38E-12 0.38E-13 0.38E-13 1
                                                                 3230.0 3225.0 3235.0 1.62E-12 3.0E-14 3.0E-14   0
FLWO/4S  B     233.73791  23.50333 4.5 48776.001234  300.0 0.05  4400.0 4200.0 4600.0 1.82E-12 1.2E-14 3.1E-14   0
FLWO/4S  V     233.73791  23.50333 4.5 48776.012012  300.0 0.05  5400.0 5200.0 5600.0 3.82E-12 1.3E-14 1.5E-14   0
FLWO/4S  R     233.73791  23.50333 4.5 48776.019013  240.0 0.05  7000.0 6200.0 7500.0 5.82E-12 1.3E-13 2.1E-13   0
FLWO/4S  I     233.73791  23.50333 4.5 48776.024988  240.0 0.08  9000.0 8200.0 9900.0 8.12E-12 3.3E-13 3.4E-13   0
\end{verbatim}
}

\clearpage

\subsection{VOTable serialization}

The VOTable version of Spectrum will represent an SED
by a series of VOTable \verb+TABLEs+. 
First there is a table with no fields
containing the overall SED and target information. Then there
is one table for each individual spectrum segment.
The data model fields described above as arrays map to
VOTable FIELDs, while the remaining fields map to PARAM.

The example below has a single RESOURCE element containing three TABLE
elements: the first TABLE element contains the top level SED object
and the second and third TABLE elements contain SED segements.

{\bf If there is only one segment in the SED, the top level SED fields
may be included in the same TABLE as the segment data }  (so that there
is only one TABLE element in all.)

{ \footnotesize
\begin{flushleft}


\begin{fmpage}

\begin{verbatim}

<?xml version="1.0" encoding="UTF-8"?>
<VOTABLE version="1.1"
  xmlns:xsi="http://www.w3.org/2001/XMLSchema-instance"
  xsi:noNamespaceSchemaLocation="http://www.ivoa.net/xml/VOTable/v1.1"
  xmlns:sed="http://www.ivoa.net/xml/SedModel/v0.94"
  xmlns="http://www.ivoa.net/xml/VOTable/v1.1">
<RESOURCE utype="sed:SED">

<TABLE utype="sed:SED">   
<PARAM name="Date" utype="sed:Date" datatype="char" arraysize="*" value="2004-05-10"/>
<PARAM name="Nseg" utype="sed:NSegments" ucd="meta.number" datatype="int" value="1"/>
<PARAM name="Type" utype="sed:Type" datatype="char" arraysize="*" value="observed"/>
<GROUP utype="sed:Target">
 <PARAM name="Target" utype="sed:Target.Name" datatype="char" arraysize="*" value="Arp 220"/>
 <PARAM name="TargetPos" utype="sed:Target.pos" unit="deg" datatype="double" arraysize="2" 
           value="233.737917 23.503330"/>
 <PARAM name="z" utype="sed:Target.redshift" datatype="float" value="0.0018"/>
</GROUP>
</TABLE>

<TABLE utype="sed:Segment">   
<!-- SegmentType can be Photometry, TimeSeries or Spectrum -->
<PARAM name="Segtype" utype="sed:Segment.SegmentType" datatype="char" 
                              arraysize="*" value="Photometry" ucd="meta.code"/>
<GROUP name="Frame" utype="sed:Segment.Frame">
 <GROUP utype="sed:Segment.Frame.Sky">
   <PARAM name="System" utype="sed:Segment.Frame.Sky.Type" ucd="frame.pos.system" 
                              datatype="char" arraysize="*" value="ICRS"/>
   <PARAM name="Equinox" utype="sed:Segment.Frame.Sky.Equinox" ucd="time.equinox;pos.eq"
                              datatype="float" value="2000.0" />
 </GROUP>
 <GROUP utype="sed:Segment.Frame.Time">
  <PARAM name="TimeFrame" utype="sed:Segment.Frame.Time.Type" ucd="frame.time.scale" 
                              datatype="char" arraysize="*" value="UTC"/>
  <PARAM name="TimeZero" utype="sed:Segment.Frame.Time.Zero" ucd="frame.time.zero" 
                              datatype="double" value="0.0"/>
 </GROUP>
 <GROUP utype="sed:Segment.Frame.SpectralCoord">
  <PARAM name="SpectralFrame" utype="sed:Segment.Frame.SpectralCoord.RefPos" ucd="frame.em.system" 
                               datatype="char" arraysize="*" value="Barycent"/>
 </GROUP>
</GROUP>
<GROUP utype="sed:Segment.Coverage">
 <GROUP utype="sed:Segment.Coverage.Location">
  <PARAM name="SkyPos" utype="sed:Segment.Coverage.Location.Sky.Value" ucd="pos.eq" unit="deg" 
                                datatype="double" arraysize="2" value="132.4210 12.1232"/>
  <PARAM name="TimeObs" utype="sed:Segment.Coverage.Location.Time.Value" ucd="time.obs" 
                                datatype="double" value="52148.3252"/>
 </GROUP>
 <GROUP utype="sed:Segment.Coverage.Extent">
  <PARAM name="SkyExtent" utype="sed:Segment.Coverage.Extent.Sky" ucd="pos.region.diameter" 
                                datatype="double" unit="arcsec" value="20"/>
  <PARAM name="TimeExtent" utype="sed:Segment.Coverage.Extent.Time" ucd="time.expo;phot.spectrum" unit="s" 
                                 datatype="double" value="1500.0" />
  <PARAM name="SpectralExtent" utype="sed:Segment.Coverage.Extent.Spectral" ucd="instr.bandwidth"
                                  unit="Angstrom" datatype="double" value="3000.0"/>
 </GROUP>

\end{verbatim}

\end{fmpage}
\begin{fmpage}

\begin{verbatim}
 <GROUP utype="sed:Segment.Coverage.Region">
  <GROUP utype="sed:Segment.Coverage.Region.Time">
   <PARAM name="TimeStart" utype="sed:Segment.Coverage.Region.Time.Start" ucd="time" unit="s" 
                                    datatype="double" value="52100.000" />
   <PARAM name="TimeStop" utype="sed:Segment.Coverage.Region.Time.Stop" ucd="time" unit="s" 
                                     datatype="double" value="52300.000" />
  </GROUP>
 </GROUP>
</GROUP>

<GROUP utype="sed:Segment.Curation">
 <PARAM name="Publisher" utype="sed:Segment.Curation.Publisher" ucd="meta.organization;meta.curation" 
                         datatype="char" arraysize="*" value="SAO"/>
 <PARAM name="PubID" utype="sed:Segment.Curation.PubID" ucd="meta.curation.pubid" 
                         datatype="char" arraysize="*" value="ivoa://cfa.harvard.edu"/>
 <PARAM name="Logo" utype="sed:Segment.Curation.Logo" ucd="meta.curation.logo" 
                         datatype="char" arraysize="*" value="http://cfa-www.harvard.edu/nvo/cfalogo.jpg"/>
 <PARAM name="Contact" utype="sed:Segment.Curation.ContactName" ucd="meta.human;meta.curation" 
                          datatype="char" arraysize="*" value="Jonathan McDowell"/>
 <PARAM name="email" utype="sed:Segment.Curation.ContactEmail" ucd="meta.email" 
                          datatype="char" arraysize="*" value="jcm@cfa.harvard.edu"/>
</GROUP>

<GROUP utype="sed:Segment.DataID">
 <PARAM name="Title" utype="sed:Segment.DataID.Title" datatype="char" arraysize="*" value="Arp 220 SED"/>
 <PARAM name="Creator" utype="sed:Segmenet.DataID.Creator" ucd="meta.curation.creator" 
                          datatype="char" arraysize="*" value="SAO/FLWO"/>
 <PARAM name="DateDate" utype="sed:Segment.DataID.Date" ucd="time;soft.dataset;meta.curation" 
                          datatype="char" arraysize="*" value="2003-12-31T14:00:02Z"/>
 <PARAM name="Version" utype="sed:Segment.DataID.Version" ucd="soft.dataset.version;meta.curation" 
                          datatype="char" arraysize="*" value="1"/>
 <PARAM name="Instrument" utype="sed:Segment.DataID.Instrument" ucd="inst.id" 
                          datatype="char" arraysize="*" value="BCS"/>
 <PARAM name="Filter" utype="sed:Segment.DataID.Collection" ucd="inst.filter.id" 
                          datatype="char" arraysize="*" value="G300"/>
 <PARAM name="CreationType" utype="sed:Segment.DataID.CreationType" 
                          datatype="char" arraysize="*" value="Archival"/>
</GROUP>

<GROUP utype="sed:Segment.Derived">
 <PARAM name="SNR" utype="sed:Segment.Derived.SNR" datatype="float" value="3.0"/>
</GROUP>

<GROUP utype="sed:Segment.Points">

<GROUP utype="sed:Segment.Points.SpectralCoord">
 <FIELDref ref="Coord"/>

 <GROUP utype="sed:Segment.Points.SpectralCoord.Accuracy">
  <FIELDref ref="BinLow"/>
  <FIELDref ref="BinHigh"/>
 </GROUP>
<!-- In this case Resolution is demoted from Field to Param since it is constant -->
 <PARAM name="Resolution" utype="sed:Segment.Points.SpectralCoord.Accuracy.Resolution" unit="Angstrom" 
                           datatype="float" value="14.2"/>
 <PARAM name="Calibration" value="ABSOLUTE"/>
</GROUP>

\end{verbatim}
\end{fmpage}


\begin{fmpage}
\begin{verbatim}
<GROUP utype="sed:Segment.Points.Flux">
 <FIELDref ref="Flux1"/>
 <GROUP utype="sed:Segment.Points.Flux.Accuracy">
  <FIELDref ref="ErrorLow"/>
  <FIELDref ref="ErrorHigh"/>
  <PARAM name="SysErr" utype="sed:Segment.Points.Flux.SysErr" unit="" datatype="float" value="0.05"/>
 </GROUP>
 <FIELDref ref="Quality"/>
</GROUP>
</GROUP>

<FIELD name="Coord" ID="Coord" utype="sed:Segment.Points.SpectralCoord.Value" ucd="em.wavelength" 
                                              datatype="double" unit="Angstrom"/>
<FIELD name="BinLow" ID="BinLow" utype="sed:Segment.Points.SpectralCoord.BinLow" 
                                    ucd="stat.min;em.wavelength" datatype="double" unit="Angstrom"/>
<FIELD name="BinHigh" ID="BinHigh" utype="sed:Segment.Points.SpectralCoord.BinHigh" 
                                    ucd="stat.max;em.wavelength" datatype="double" unit="Angstrom"/>
<FIELD name="Flux" ID="Flux1" utype="sed:Segment.Points.Flux.Value" ucd="phot.flux;em.wavelength" 
                                    datatype="double" unit="erg cm**(-2) s**(-1) Angstrom**(-1)"/>
<FIELD name="ErrorLow" ID="ErrorLow" utype="sed:Segment.Points.Flux.Accuracy.StatErrLow" 
                                    datatype="double" unit="erg cm**(-2) s**(-1) Angstrom**(-1)"/>
<FIELD name="ErrorHigh" ID="ErrorHigh" utype="sed:Segment.Points.Flux.Accuracy.StatErrHigh" 
                                    datatype="double" unit="erg cm**(-2) s**(-1) Angstrom**(-1)"/>
<FIELD name="Quality" ID="Quality" datatype="int" utype="sed:Segment.Points.Flux.Quality"/>

<DATA>
<TABLEDATA>
<!-- Note slightly nonlinear wavelength solution -->
<!-- Second row is upper limit -->
<!-- Third row has quality mask set -->
<TR><TD>3200.0</TD><TD>3195.0</TD><TD>3205.0</TD><TD>1.38E-12</TD><TD>5.2E-14</TD> 
                                                   <TD>6.2E-14</TD><TD>0</TD></TR>
<TR><TD>3210.5</TD><TD>3205.0</TD><TD>3216.0</TD><TD>1.12E-12</TD><TD>1.12E-12</TD>
                                                    <TD>0</TD><TD>0</TD></TR>
<TR><TD>3222.0</TD><TD>3216.0</TD><TD>3228.0</TD><TD>1.42E-12</TD><TD>1.3E-14</TD>
                                                    <TD>0.2E-14</TD><TD>3</TD></TR>
</TABLEDATA>
</DATA>
</TABLE>

<!-- The second table below gives a radio light curve, and uses more defaults. -->
<TABLE utype="sed:Segment">   

<PARAM name="Segtype" utype="sed:Segment.SegmentType" datatype="char" arraysize="*" 
                    value="TimeSeries" ucd="meta.code"/>
<GROUP name="Frame" utype="sed:Segment.Frame">
 <GROUP utype="sed:Segment.Frame.Sky">
   <PARAM name="System" utype="sed:Segment.Frame.Sky.Type" ucd="frame.pos.system" 
                               datatype="char" arraysize="*" value="FK5"/>
   <PARAM name="Equinox" utype="sed:Segment.Frame.Sky.Equinox" ucd="time.equinox;pos.eq" 
                               datatype="double" value="2000.0"/>
 </GROUP>
 <GROUP utype="sed:Segment.Frame.Time">
  <PARAM name="TimeType" utype="sed:Segment.Frame.Time.Type" datatype="char" arraysize="*" value="MJD"/>
  <PARAM name="TimeZero" utype="sed:Segment.Frame.Time.Zero" ucd="frame.time.zero" datatype="double" value="0.0"/>
  <PARAM name="TimeSIDim" utype="sed:Segment.Frame.Time.SIDim" datatype="char" arraysize="*" value="T"/>
  <PARAM name="TimeSystem"  utype="sed:Segment.Frame.Time.System" ucd="frame.time.scale" 
                                datatype="char" arraysize="*" value="TT"/>
  <PARAM name="TimeRefPos"  utype="sed:Segment.Frame.Time.RefPos" ucd="pos.eq.ra" 
                                datatype="char" arraysize="*" value="TOPOCENTER"/>
 </GROUP>
\end{verbatim}
\end{fmpage}

\begin{fmpage}
\begin{verbatim}
 <GROUP utype="sed:Segment.Frame.SpectralCoord">
  <PARAM name="SpectralRefPos" utype="sed:Segment.Frame.SpectralCoord.RefPos" ucd="frame.em.system" 
                                datatype="char" arraysize="*" value="Barycent"/>
  <PARAM name="SpecSIDim"  utype="sed:Segment.Frame.SpectralCoord.SIDim" 
                                datatype="char" arraysize="*" value="10-10 L"/>
 </GROUP>
 <GROUP utype="sed:Segment.Frame.Flux">
  <PARAM name="FluxSIDim"  utype="sed:Segment.Frame.Flux.SIDim" 
                                datatype="char" arraysize="*" value="10+7 ML-1T-3"/>
 </GROUP>
</GROUP>
<GROUP utype="sed:Segment.Coverage">
 <GROUP utype="sed:Segment.Coverage.Location">
  <PARAM name="SkyPos" utype="sed:Segment.Coverage.Location.Sky.Value" ucd="pos.eq" unit="deg" 
                                datatype="double" arraysize="2" value="132.4210 12.1231"/>
  <PARAM name="TimeObs" utype="sed:Segment.Coverage.Location.Time.Value" ucd="time.obs" 
                                datatype="double" value="0.0"/>
  <PARAM name="SpectralValue" utype="sed:Segment.Coverage.Location.SpectralCoord.Value" 
                                ucd="em.freq" unit="MHz" datatype="double" value="1660.0"/>
 </GROUP>
 <GROUP utype="sed:Segment.Coverage.Extent">
  <PARAM name="SkyExtent" utype="sed:Segment.Coverage.Extent.Sky" ucd="pos.region.diameter" 
                                unit="arcsec" datatype="double" value="2.5"/>
  <PARAM name="TimeExtent" utype="sed:Segment.Coverage.Extent.Time" ucd="time" unit="d" 
                                datatype="double" value="1095.0"/>
  <PARAM name="SpectralExtent" utype="sed:Segment.Coverage.Extent.Spectral" ucd="instr.bandwidth" 
                                unit="MHz" datatype="double" value="400"/>
 </GROUP>
</GROUP>

<GROUP utype="sed:Curation">
 <PARAM name="Publisher" utype="sed:Segment.Curation.Publisher" ucd="meta.organization;meta.curation"
                                          datatype="double" arraysize="*" value="NRAO"/>
</GROUP>

<GROUP utype="sed:DataID">
 <PARAM name="Title" utype="sed:Segment.DataID.Title" datatype="char" arraysize="*" 
                                       value="Arp 220 1660 MHz Monitoring"/>
 <PARAM name="Creator" utype="sed:Segment.DataID.Creator" ucd="meta.curation.creator" 
                                       datatype="char" arraysize="*" value="NRAO"/>
 <PARAM name="Instrument" utype="sed:Segment.DataID.Instrument" ucd="inst.id"
                                      datatype="char"  arraysize="*" value="VLA"/>
</GROUP>

<GROUP utype="sed:Segment.Points.Time">
 <FIELDref ref="Time"/>
</GROUP>

<GROUP utype="sed:Segment.Points.SpectralCoord">
 <PARAM name="Coord" utype="sed:Segment.Points.SpectralCoord.Value" ucd="em.freq" unit="MHz" 
                                               datatype="double" value="1660.0"/>
</GROUP>

\end{verbatim}
\end{fmpage}

\begin{fmpage}
\begin{verbatim}
<GROUP utype="sed:Segment.Points.Flux">
 <FIELDref ref="Flux2"/>
 <GROUP utype="sed:Segment.Points.Flux.Accuracy">
  <PARAM name="ErrorLow" utype="sed:Segment.Points.Flux.Accuracy.StatErrLow" unit="mJy" 
                                                  datatype="double" value="0.15"/>
  <PARAM name="ErrorHigh" utype="sed:Segment.Points.Flux.Accuracy.StatErrHigh" unit="mJy" 
                                                   datatype="double" value="0.15"/>
  <PARAM name="SysErr" utype="sed:Segment.Points.Flux.Accuracy.SysErr" datatype="double" unit="" value="0.10"/>
 </GROUP>
</GROUP>

<FIELD name="Time" ID="Time" utype="sed:Segment.Points.Time.Value" ucd="time" datatype="float" unit="d"/>
<FIELD name="Flux" ID="Flux2" utype="sed:Segment.Points.Flux.Value" ucd="phot.flux;em.freq" 
                                                         datatype="double" unit="mJy"/>

<DATA>
<TABLEDATA>
<TR><TD>46066.5</TD><TD>5.18</TD></TR>
<TR><TD>46246.5</TD><TD>6.21</TD></TR>
<TR><TD>46431.0</TD><TD>11.32</TD></TR>
<TR><TD>46796.0</TD><TD>15.00</TD></TR>
<TR><TD>46975.3</TD><TD>7.32</TD></TR>
<TR><TD>47161.5</TD><TD>3.10</TD></TR>
</TABLEDATA>
</DATA>
</TABLE>

</RESOURCE>
</VOTABLE>
\end{verbatim}
\end{fmpage}

\end{flushleft}
}

\clearpage
\subsection{XML serialization: photometry example}

{ \footnotesize
\begin{flushleft}


\begin{fmpage}



%\colorbox{iblue}{
\begin{verbatim}
<?xml version="1.0" encoding="UTF-8"?>
<SED xmlns:xsd="http://www.w3.org/2001/XMLSchema" xmlns:xsi="http://www.w3.org/2001/XMLSchema-instance" 
xmlns="http://ivoa.net/xml/SED/SED-0.96.xsd">
<!-- xml instance example -->
 <Date>2004-05-10</Date>
 <NSegments ucd="meta.number">2</NSegments>
 <Target>
  <name>Arp 220 </name>
  <pos>233.737917 23.503330</pos>
  <redshift>0.0018</redshift>  
 </Target>
 <Segment>
  <SegmentType ucd="meta.code">Spectrum</SegmentType>
  <Frame>
   <Sky>
    <Equinox ucd="time.equinox;pos.eq">2000.0</Equinox>
    <System ucd="frame.pos.system">ICRS</System>
   </Sky>
   <Time>
     <TimeType ucd="frame.time.system?">MJD</TimeType> 
     <TimeZero ucd="frame.time.zero">0.0 </TimeZero>
     <TimeSystem ucd="frame.time.scale">UTC</TimeFrame> 
     <TimeRefPos ucd="pos;frame.time">TOPOCENTER</TimeZero>
   </Time>
   <SpectralCoord>
     <SpectralRefPos ucd="frame.em.system">Barycent</SpectralRefPos>
   </SpectralCoord>
  </Frame>
  <Coverage>
  <Location>
    <Sky ucd="pos.eq" unit="deg">132.4210 12.1232</Sky>
    <Time ucd="time.obs" unit="d">52148.3252</Time>
   </Location>
   <Extent>
     <Sky ucd="pos.region.diameter" unit="arcsec">20</Sky>
     <Time ucd="time.expo;phot.spectrum" unit="s">1500.0</Time>
     <Spectral ucd="instr.bandwidth" unit="Angstrom">3000.0</Spectral>
   </Extent>
  </Coverage>
  <Curation>
   <Publisher ucd="meta.organization;meta.curation">SAO</Publisher>
   <PubID ucd="meta.curation.pubid">ivoa://cfa.harvard.edu</PubID>
   <Logo ucd=meta.curation.logo">http://cfa-www.harvard.edu/nvo/cfalogo.jpg</Logo>
   <Contact>
     <Name ucd="meta.human;meta.curation">Jonathan McDowell</Name>
     <Email ucd="meta.email">jcm@cfa.harvard.edu</Email>
   </Contact>
  </Curation>
  <DataID>
    <Title>Arp 220 SED</Title>
    <Creator ucd="meta.curation.creator">SAO/FLWO</Creator>
    <Date ucd="time;soft.dataset;meta.curation">2003-12-31T14:00:02Z</Date>
    <Version ucd="soft.dataset.version;meta.curation">1</Version>
    <Instrument ucd="inst.id">BCS</Instrument>
    <Collection>Archival</Collection>
  </DataID>
  <Derived>
    <SNR>3.0</SNR>
  </Derived>
\end{verbatim}
\end{fmpage}

\begin{fmpage}
\begin{verbatim}

  <!-- Define table structure -->
  <Fields>
    <SpectralCoord>
     <Value ucd="em.wavelength" unit="Angstrom"/>    
     <Accuracy>
      <BinLow ucd="stat.min;em.wavelength" unit="Angstrom"/> 
      <BinHigh ucd="stat.max;em.wavelength" unit="Angstrom"/>
     </Accuracy>
   </SpectralCoord>
   <Flux>
     <Value ucd="phot.flux;em.wavelength" unit="erg cm**(-2) s**(-1) Angstrom**(-1)"/>
     <Accuracy>
       <StatErrLow  unit="erg cm**(-2) s**(-1) Angstrom**(-1)"/>
       <StatErrHigh unit="erg cm**(-2) s**(-1) Angstrom**(-1)"/>
       <SysErr>0.05</SysErr>
     </Accuracy>
     </Quality>
   </Flux>
  </Fields>
  <!-- Use table structure -->
  <Points>
   <Point> 
    <SpectralCoord>
     <Value>3200.0</Value>
     <Accuracy><BinLow>3195.0</BinLow><BinHigh>3205.0</BinHigh></Accuracy>
    </SpectralCoord>
    <Flux>
      <Value>1.38E-12</Value>
      <Accuracy><StatErrLow>5.2E-14</StatErrLow><StatErrHigh>6.2E-14</StatErrHigh></Accuracy>
      <Quality>0</Quality>
    </Flux>
   </Point>

   <Point> 
    <SpectralCoord>
     <Value>3210.5</Value>
     <Accuracy><BinLow>3205.0</BinLow><BinHigh>3216.0</BinHigh></Accuracy>
    </SpectralCoord>
    <Flux>
      <Value>1.12E-12</Value>
      <Accuracy><StatErrLow>1.12E-12</StatErrLow><StatErrHigh>0</StatErrHigh></Accuracy>
      <Quality>0</Quality>
    </Flux>
   </Point>

   <Point> 
    <SpectralCoord>
     <Value>3222.0</Value>
     <Accuracy><BinLow>3216.0</BinLow><BinHigh>3228.0</BinHigh></Accuracy>
    </SpectralCoord>
    <Flux>
      <Value>1.42E-12</Value>
      <Accuracy><StatErrLow>1.3E-14</StatErrLow><StatErrHigh>0.2E-14</StatErrHigh></Accuracy>
      <Quality>3</Quality>
    </Flux>
   </Point>

  </Points>
 </Segment>

\end{verbatim}
\end{fmpage}

\begin{fmpage}
\begin{verbatim}


 <Segment>
  <SegmentType ucd="meta.code">TimeSeries</SegmentType>
  <Frame>
   <Sky>
    <Equinox ucd="time.equinox;pos.eq">2000.0</Equinox>
    <System ucd="frame.pos.system">FK5</System>
   </Sky>
   <Time>
     <TimeType ucd="frame.time.type">MJD</TimeType> 
     <TimeFrame ucd="frame.time.scale">TT</TimeFrame> 
     <TimeZero ucd="frame.time.zero">0.0 </TimeZero>
   </Time>
   <SpectralCoord>
     <SpectralRefPos ucd="frame.em.system">Barycent</SpectralRefPos>
   </SpectralCoord>
  </Frame>

  <Coverage>
   <Location>
    <Sky ucd="pos.eq" unit="deg">132.4210 12.1231</Sky>
    <Time ucd="time.obs" unit="d">0</Time>
   </Location>
   <Extent>
     <Sky ucd="pos.region.diameter" unit="arcsec">2.5</Sky>
     <Time ucd="time.expo;phot.spectrum" unit="s">1095.0</Time>
     <Spectral ucd="instr.bandwidth" unit="Angstrom">400.0</Spectral>
   </Extent>
  </Coverage>

  <Curation>
   <Publisher ucd="meta.organization;meta.curation">NRAO</Publisher>
  </Curation>

  <DataID>
    <Title>Arp 220 1660 MHz Monitoring</Title>
    <Creator ucd="meta.curation.creator">NRAO</Creator>
    <Instrument ucd="inst.id">VLA</Instrument>
  </DataID>

  <Fields>
   <Time>
     <Value ucd="time" unit="d"/>    
   </Time>
   <SpectralCoord>
     <Value ucd="em.freq" unit="MHz"> 1660.0   </Value>    
   </SpectralCoord>
   <Flux>
     <Value ucd="phot.flux;em.freq" unit="mJy"/>
     <Accuracy>
       <StatErrLow  unit="mJy">0.15 </StatErrLow>
       <StatErrHigh unit="mJy">0.15 </StatErrHigh>
       <SysErr>0.10</SysErr>
     </Accuracy>
   </Flux>
  </Fields>

\end{verbatim}
\end{fmpage}

\begin{fmpage}
\begin{verbatim}
  <Points>

   <Point> 
    <Time>
     <Value>46066.5</Value>
    </Time>
    <Flux>
      <Value>5.18</Value>
    </Flux>
   </Point>

   <Point> 
    <Time>
     <Value>46246.5</Value>
    </Time>
    <Flux>
      <Value>6.21</Value>
    </Flux>
   </Point>
  <!-- etc. -->
  </Points>
</Segment>

</SED>

\end{verbatim}
%}
\end{fmpage}
\end{flushleft}
}


\section*{References}


\vskip 0.1in
\par\noindent
Hanisch, R., (ed)., Resource Metadata for the VO, Version 1.01,
2004 Apr 26.  \\
http://www.ivoa.net/Documents/latest/RM.html

\vskip 0.1in
\par\noindent
Derriere, S. et al (eds.), UCD, Moving to UCD 1+, 2004 Oct 26.
\\
http://www.ivoa.net/Documents/latest/UCD.html

\end{document}
